\chapterhead{5}{Groups Inside Groups}{Sorting real animals into smaller and smaller circles}\label{ch5}

Real animals are sorted the same way as bugs, with groups inside groups. All the animals with six legs form a group. Inside that group are smaller groups, such as ones with big colourful wings. Here is a key for eight animals.

\begin{center}
\sffamily\small\renewcommand{\arraystretch}{1.5}
\begin{tabular}{@{}p{0.6cm}p{6.6cm}p{2.9cm}@{}}
\toprule
\textbf{No.} & \textbf{Question} & \textbf{Yes, then / No, then}\\
\midrule
1 & Does it have six legs? & go to 2 / go to 3\\
2 & Does it have large, colourful wings? & butterfly / ant\\
3 & Does it have eight legs? & spider / go to 4\\
4 & Does it have feathers or fur? & go to 5 / go to 6\\
5 & Does it have feathers? & sparrow / cat\\
6 & Does it have a hard shell on its back? & snail / go to 7\\
7 & Does it have fins? & goldfish / frog\\
\bottomrule
\end{tabular}
\end{center}

\begin{center}
\begin{tikzpicture}
  \node[bigcir,minimum size=4.2cm] (a) at (0,0) {};
  \node[font=\sffamily\scriptsize] at (0,1.7) {8 animals};
  \node[cir,minimum size=1.6cm,font=\sffamily\tiny,align=center] at (-1.0,0.3) {six legs\\butterfly, ant};
  \node[cir,minimum size=1.5cm,font=\sffamily\tiny,align=center] at (1.0,0.7) {eight legs\\spider};
  \node[cir,minimum size=1.7cm,font=\sffamily\tiny,align=center] at (1.0,-0.8) {fur or feathers\\sparrow, cat};
  \node[cir,minimum size=1.5cm,font=\sffamily\tiny,align=center] at (-1.0,-1.1) {shell, fins\\snail, fish, frog};
  \node[note,anchor=west,text width=4.6cm,align=left] at (3.0,0) {Each smaller circle sits inside the big one. A good key makes circles that do not overlap.};
\end{tikzpicture}
\end{center}

\begin{tryit}{Try the key on a walk}
\textbf{You need:} paper and a pencil, or a camera, and a safe place to look.
\begin{enumerate}[leftmargin=1.6em]
\item Look for three or more animals outside, from a safe distance. Draw or photograph each. Do not touch or pick up any animal.
\item For each, follow the key from question 1. Write where it ends.
\item If an animal does not fit, write down what question it broke.
\end{enumerate}
\par\medskip
\renewcommand{\arraystretch}{1.8}
\begin{tabular}{|l|c|c|}
\hline
\textbf{Animal I saw} & \textbf{Where the key ended} & \textbf{Did it fit?}\\ \hline
 & & \\ \hline
 & & \\ \hline
 & & \\ \hline
\end{tabular}
\end{tryit}
\adult{Look from a safe distance. Do not touch plants or animals you cannot name, and wash hands afterwards.}

\begin{worked}{Practice example}
A sparrow: Q1 six legs? no. Q3 eight legs? no. Q4 feathers or fur? yes. Q5 feathers? yes. It ends at sparrow, after 4 questions.\\
A frog: Q1 no, Q3 no, Q4 no, Q6 no, Q7 no. It ends at frog, after 5 questions.
\end{worked}

\begin{multix}[hmm]\multixbot[hmm]\tcblower
A dog has fur and no feathers, so the key puts it in the cat group. That means the key calls a dog a cat. I think the key is broken. Is it?
\end{multix}

\trust{We followed the key for two animals, counted the questions and checked each ended in a group of its own.}
{The sparrow ended in 4 questions and the frog in 5, and each ended in a different place.}
{That a dog or a worm would end somewhere sensible. The key only knows the eight animals it was built for.}

\begin{keepbox}
We used a key for eight animals. \textbf{We kept} \wline[3.4cm] \textbf{and we set aside} \wline[3.4cm].
\end{keepbox}

\clearpage

\begin{puzzle}{An ant}
Follow the key for an ant. Write the question numbers you visit. \wline[5cm]
\end{puzzle}

\begin{puzzle}{A frog}
How many questions does the frog need? \wline[1.6cm]
\end{puzzle}

\begin{puzzle}{A goldfish}
Write the question numbers for a goldfish. \wline[5cm]
\end{puzzle}

\begin{puzzle}{The fewest}
Which animal needs the fewest questions? \wline[3cm] How many? \wline[1.6cm]
\end{puzzle}

\begin{puzzle}{Counting a group}
How many of the eight animals have six legs? \wline[1.6cm] How many do not? \wline[1.6cm]
\end{puzzle}

\clearpage

\begin{puzzle}{A dog}
Where does the key end for a dog? \wline[3cm] What does that tell you about the ``cat'' group? \wline[5cm]
\end{puzzle}

\begin{puzzle}{Spiders and insects}
Insects have six legs. Use the key to say whether a spider is an insect. \wline[3cm] Why? \wline[4cm]
\end{puzzle}

\begin{puzzle}{Another way to sort}
Sort the eight animals by where they live instead: water, land or air. Which feature did you use? \wline[3cm] Is it the same as the key's groups? \wline[3cm]
\end{puzzle}

\begin{puzzle}{Your walk}
List three animals from your walk and where the key put each. \wline[6cm]
\end{puzzle}

\begin{puzzle}{Draw the circles}
Eight animals are split by six legs into 2 and 6. Fill in the circles.
\begin{center}
\begin{tikzpicture}
  \node[cir] (a) at (0,0) {8}; \node[cir,dashed] (b) at (4,0.9) {}; \node[cir,dashed] (c) at (4,-0.9) {};
  \draw[popA] (a) -- (b); \draw[popA] (a) -- (c);
\end{tikzpicture}
\end{center}
\end{puzzle}
\clearpage
\thispagestyle{fancy}
\subhead{Animal log}
Record the animals you see, and follow the key for each.\par\medskip
\renewcommand{\arraystretch}{3.0}
\begin{tabularx}{\linewidth}{|>{\sffamily\small\bfseries}p{2.4cm}|X|X|X|}
\hline
Animal & Where I saw it & Where the key ended & Did it fit? \\ \hline
 & & & \\ \hline
 & & & \\ \hline
 & & & \\ \hline
 & & & \\ \hline
 & & & \\ \hline
\end{tabularx}

\vspace{8pt}
Which animal broke the key, if any? \wline[6cm]\\[10pt]
Draw one animal you saw.\par\medskip
\begin{tikzpicture}\draw[ink!70,line width=1pt,rounded corners=3pt] (0,0) rectangle (11.4,4.2);\end{tikzpicture}
\clearpage
